\section{Implementation}
\label{sec:implementation}

This section summarizes the FeBEx implementation at a high level, focusing on packet-processing flow rather than code-level details.

\subsection{Components}

FeBEx is implemented with three components:
\begin{itemize}
    \item \textbf{P4 data plane} (BMv2, v1model): parses FeBEx metadata, performs tenant steering, deduplication, and receipt cloning.
    \item \textbf{P4Runtime controller} (Python + finsy): installs LPM tenant rules, configures clone session, and rotates epochs.
    \item \textbf{Experiment harness} (Mininet + Scapy): generates traffic, captures outputs, and computes metrics.
\end{itemize}

\subsection{Switch Pipeline Flow}

\begin{figure}[H]
\centering
\fbox{\parbox{0.94\linewidth}{
\textbf{Switch Ingress Flow (high-level)}\\[1mm]
\textit{Parse packet} (Ethernet/IPv4/UDP/FeBExMeta)\\
$\downarrow$\\
\textit{Tenant steering} (LPM on DevAddr $\rightarrow$ select tenant egress port)\\
$\downarrow$\\
\textit{Dedup check} (hash index + key, compare with register state + epoch)\\
$\downarrow$\\
\textbf{Duplicate?}\\
No $\rightarrow$ update register, forward to tenant, clone receipt to cloud\\
Yes $\rightarrow$ drop packet
}}
\caption{FeBEx switch pipeline at implementation level.}
\label{fig:impl-switch-flow}
\end{figure}

\subsection{Controller Flow}

\begin{figure}[H]
\centering
\fbox{\parbox{0.94\linewidth}{
\textbf{Controller Lifecycle (high-level)}\\[1mm]
Connect via P4Runtime\\
$\downarrow$\\
Load P4 pipeline + install tenant LPM entries\\
$\downarrow$\\
Set dedup mode and register defaults\\
$\downarrow$\\
Configure clone session (session 100 $\rightarrow$ cloud port)\\
$\downarrow$\\
Run periodic epoch rotation thread (every configured interval)
}}
\caption{Controller behavior used in all experiments.}
\label{fig:impl-controller-flow}
\end{figure}

\subsection{Deduplication State Design}

FeBEx keeps only compact state in registers:
\begin{itemize}
    \item \texttt{dedup\_keys[index]}: 32-bit verification key.
    \item \texttt{dedup\_epochs[index]}: 16-bit epoch when slot was written.
    \item \texttt{current\_epoch[0]}: controller-rotated epoch value.
    \item \texttt{dedup\_enabled[0]}: runtime A/B toggle for ON/OFF comparisons.
\end{itemize}

This design avoids expensive per-device state while preserving line-rate-friendly lookup and update operations.

\subsection{Switch Variants for Boundary-Leakage Mitigation}

To study boundary effects and mitigation tradeoffs, we implement and compare three switch variants:

\begin{itemize}
    \item \textbf{V1 (baseline)}: single-epoch check using \texttt{stored\_epoch == current\_epoch}. This is vulnerable to boundary leakage when duplicates straddle an epoch rotation.
    \item \textbf{V2 (sliding window)}: checks both current and previous epoch (logical two-epoch window), so cross-boundary duplicates are still recognized and suppressed.
    \item \textbf{V3 (dual-register Bloom)}: uses a dual-register Bloom-style design to reduce hash-collision false positives (stronger collision behavior), but its boundary behavior remains similar to V1.
\end{itemize}

\begin{figure}[H]
\centering
\fbox{\parbox{0.96\linewidth}{
\textbf{Boundary-Leakage Mitigation Across Variants}\\[1mm]
\textbf{Timeline:} \hspace{1mm} ... \(E_t\) \(\mid\) \(E_{t+1}\) ...  (epoch boundary at \(\mid\))\\[1mm]
\textbf{V1:} check \(E_{t+1}\) only \\
\hspace*{7mm} duplicate that first arrived in \(E_t\) and repeats in \(E_{t+1}\) may leak\\[1mm]
\textbf{V2 (sliding-window):} check \(E_t\) \textit{or} \(E_{t+1}\) \\
\hspace*{7mm} cross-boundary duplicates are still recognized and suppressed\\[1mm]
\textbf{V3 (dual-register Bloom):} collision-focused design with stronger FP resistance \\
\hspace*{7mm} improves collision behavior, but does not eliminate boundary leakage by itself
}}
\caption{Conceptual comparison of V1, V2 (sliding window), and V3 (dual-register Bloom).}
\label{fig:v123-boundary-mitigation}
\end{figure}

\subsection{Packet and Traffic Path}

Each gateway sends UDP packets with FeBEx metadata: $(\texttt{dev\_addr},\texttt{fcnt},\texttt{gw\_id})$. The switch forwards exactly one copy of each unique uplink to the correct tenant backend and mirrors that first-forwarded copy to cloud for Proof-of-Coverage attribution.

\subsection{Implementation Choices}

\begin{itemize}
    \item \textbf{Two-hash approach}: one hash for index, one for verification key to reduce collision ambiguity.
    \item \textbf{Epoch expiry}: no explicit deletion; stale entries naturally age out when epoch changes.
    \item \textbf{LPM tenant steering}: simple multi-tenant isolation without per-device forwarding entries.
    \item \textbf{Clone-on-first-forward}: preserves attribution while still suppressing redundant backhaul packets.
\end{itemize}
